% =====================================================================
%  preamble.tex —— 中文版《Knot Invariants from Four-Dimensional Gauge Theory》导言区
%  编译方式： XeLaTeX（推荐，ctex 自动处理中文字体）
%  英文原稿： D. Gaiotto & E. Witten, arXiv:1106.4789v1, 120 pp.
%  说明： 原稿为 JHEP 风格（jheppub），此处改为 article+ctex，
%         并保留原稿全部自定义数学宏，以确保公式逐字无损。
% =====================================================================

% ---------- 版式与字体 ----------
\usepackage[UTF8]{ctex}
\usepackage[margin=2.4cm]{geometry}

% ---------- 数学 ----------
\usepackage{amsmath}
\usepackage{amssymb}
\usepackage{amsthm}
\usepackage{bm}

% ---------- 图形与链接 ----------
\usepackage{graphicx}
\graphicspath{{figs/}}
\usepackage{xcolor}
\usepackage{longtable}
\usepackage{booktabs}
\usepackage[most]{tcolorbox}
\usepackage{hyperref}
\hypersetup{
  colorlinks=true,
  linkcolor=blue!45!black,
  citecolor=green!35!black,
  urlcolor=blue!55!black,
  pdftitle={四维规范场论中的纽结不变量（中译本）},
  pdfauthor={Davide Gaiotto, Edward Witten}
}

% =====================================================================
%  原稿自定义数学宏（逐字保留；仅将 TeX 传统字体族 cmmib/eurm/eusm
%  换为 XeLaTeX 下等价的外观：粗体数学斜体 / 罗马 / 花体）
% =====================================================================
\def\t{{ \sf t}}
\def\ct{{\cmmib t}}
\def\s{{s}}
\def\be{\begin{equation}}
\def\ee{\end{equation}}
\def\bt{{{t}}}
\def\hat{\widehat}
\def\tilde{\widetilde}
\def\frak{\mathfrak}
\def\D{{\mathcal D}}
\def\MM{M_4}
\def\k{{\cmmib k}}
\def\Bbb{\mathbb}
\def\Mthree{M_3}
\def\A{{\mathcal A}}
\def\d{{\mathrm d}}
\def\tW{{W}}
\def\R{{\mathbb R}}
\def\C{{\mathbb C}}
\def\U{{\mathcal U}}
\def\CP{{\mathbb{CP}}}
\def\N{{\mathcal N}}
\def\T{{\mathcal T}}
\def\F{{\mathcal F}}
\def\Tr{{\mathrm {Tr}}}
\def\Z{{\mathbb Z}}
\def\Q{{\mathbb Q}}
\def\CC{{\mathcal C}}
\def\ad{{\mathrm{ad}}}
\def\Btriv{{\mathcal B}_{\mathrm{triv}}}
\def\L{{ L}}
\def\Bcc{{\mathcal B_{\mathrm{cc}}}}
\def\K{{\mathcal K}}
\def\V{{\mathcal V}}
\def\J{{\mathcal J}}
\def\I{{\mathcal I}}
\def\P{{\mathcal P}}
\def\B{{\mathcal B}}
\def\Boper{{\mathcal B_{\mathrm{oper}}}}
\def\NS5{\mathrm{NS5}}
\def\uD5{\mathrm{D5}}
\def\M{{\mathcal M}}
\def\W{{\mathcal W}}
\def\H{{\mathcal H}}
\def\ZZ{\eusm B}
\def\cZ{\mathcal Z}
\def\sW{{ W}}
\def\sY{{ Y}}
\def\nY{{\mathcal Y}}
\def\ca{{\cmmib a}}
\def\ss{{d}}
\def\m{\cmmib m}
\def\bar{\overline}
% ---- 字体族替代（XeLaTeX 现代等价物） ----
\newcommand{\cmmib}[1]{\boldsymbol{#1}}      % 原 cmmib10：粗体数学斜体
\newcommand{\eusm}[1]{\mathcal{#1}}          % 原 eusm10：欧拉花体
\newcommand{\eurm}[1]{\mathrm{#1}}           % 原 eurm10：欧拉罗马

% =====================================================================
%  结构化学术版式
% =====================================================================

% ---- “本节重点”方框 ----
\newtcolorbox{keytakeaways}{
  enhanced,
  breakable,
  colback=blue!3!white,
  colframe=blue!45!black,
  boxrule=0.7pt,
  arc=2pt,
  left=6pt, right=6pt, top=5pt, bottom=5pt,
  title={\bfseries 本节重点（Key Takeaways）},
  fonttitle=\small\bfseries,
  coltitle=white,
  attach boxed title to top left={yshift=-2mm, xshift=4mm},
  boxed title style={colback=blue!45!black, arc=1pt, boxrule=0pt},
  top=10pt
}

% ---- 译者按 ----
\newtcolorbox{transnote}{
  enhanced,
  breakable,
  colback=gray!4!white,
  colframe=gray!55!black,
  boxrule=0.5pt,
  arc=1pt,
  left=5pt, right=5pt, top=4pt, bottom=4pt,
  title={\bfseries 译者按},
  fonttitle=\small,
  coltitle=black,
  colbacktitle=gray!20!white,
  top=9pt
}

% ---- 行距与段落 ----
\linespread{1.08}
\setlength{\parskip}{2pt}

% ---- 图表标题 ----
\usepackage[font=small,labelfont=bf]{caption}
